\begin{table}[t]
\centering\small
\setlength{\tabcolsep}{5pt}
\begin{tabular}{@{}llrrr@{}}
\toprule
& & \textbf{Full set} & \textbf{Context-complete} & \textbf{Context-missing} \\
\midrule
\multicolumn{2}{@{}l}{\textbf{All questions}} & \textbf{10{,}198} & \textbf{7{,}625} & \textbf{2{,}573} \\
\midrule
\multirow{2}{*}{Access}
 & Public      & 5{,}110 & 3{,}406 & 1{,}704 \\
 & Held out    & 5{,}088 & 4{,}219 &   869 \\
\midrule
\multirow{3}{*}{Difficulty}
 & Easy        & 6{,}578 & 5{,}929 &   649 \\
 & Medium      & 2{,}183 & 1{,}324 &   859 \\
 & Hard        & 1{,}437 &   372 & 1{,}065 \\
\bottomrule
\end{tabular}
\caption{The two evaluation sets. \texttt{difficulty\_v1} is a single classification applied to
all $10{,}198$ items, and the context-complete set is a filter over items rather than a second
classification, so no item changes band between the two columns. An item is context-missing when
it cannot be answered from its own stem and options, which happens when a case vignette or an
exhibit table was delivered as an image the collection pipeline did not capture. Every released
record carries both its band and its \texttt{context\_status}, so either column is reproducible
from the shipped files.}
\label{tab:sets}
\end{table}
